\chapter*{Abstract}
\addcontentsline{toc}{chapter}{Abstract} 

Private shuttle transportation services in Uganda, particularly on intercity routes such as Kampala–Kabale, face significant challenges related to coordination, booking management, and operational efficiency. These services largely rely on manual processes, including phone-based bookings and informal passenger coordination, which result in inefficiencies such as unreliable seat allocation, unutilized capacity, and limited visibility for both operators and passengers.

This study addressed these challenges through the design and development of Traxis, a mobile application aimed at improving the organization and efficiency of private shuttle operations. The system was designed to support seat-based booking, real-time vehicle and occupancy tracking, and structured scheduling of shuttle services. It also incorporated features for managing passenger information, monitoring vehicle movement, and facilitating communication between operators and users.

Traxis adopted a dual-interface approach, providing functionality for both passengers and operators. Passengers were able to search for available routes, book specific seats, and receive real-time updates, while operators were able to monitor trips, track vehicle activity, and manage bookings through an administrative dashboard. The system design emphasized usability and accessibility to accommodate users with varying levels of digital literacy.

The project demonstrated the potential of mobile technology to improve coordination and resource utilization within privately operated shuttle services. By introducing structured booking, real-time tracking, and integrated operational management, Traxis provided a scalable and context-aware solution for enhancing efficiency, transparency, and reliability in Uganda’s intercity transport sector.